% Appendix section of theory Part 0 (input from main.tex). Companion to Definition def:cycle-eps.
\section{Appendix: why the product-kernel defect does not imply Lemma~\ref{lem:eps}}
\label{app:counterexample}

Two cycle-level definitions of the controller defect compete. (P)~The
\emph{product-kernel defect}
$\epsbar{ctrl}^{\mathrm{prod}} := \sup_{x_{0:\Tc-1}} \dtv\bigl(\bigotimes_t
P_{\mathrm{ctrl}}(\cdot\mid x_{t-1},\theta_t),\ \bigotimes_t
P^{\sigma}_{\mathrm{ctrl}}(\cdot\mid x_{t-1},\theta_t)\bigr)$ --- the
equivariance defect of the $\Tc$-fold product kernel with the conditioning
trajectory held \emph{fixed}; (H)~the \emph{closed-loop hybrid displacement}
of \eqref{eq:cycle-eps}. Only (H) bounds the left-hand side of
\eqref{eq:eps-bound}: (P) can be strictly smaller than the closed-loop
displacement, so a lemma stated in (P) would be false.

\paragraph{Setting.}
$\Tc=2$, binary inputs $u_1,u_2\in\{0,1\}$, perfect-memory dynamics
$x_1=u_1$ ($x_0$ fixed), and, with $0<\varepsilon\le\eta<\tfrac12$, nominal
and conjugated controller kernels
$A_1=\mathrm{Bern}(\tfrac12)$, $A_1'=\mathrm{Bern}(\tfrac12+\varepsilon)$,
$A_2(\cdot\mid x)=\mathrm{Bern}(p_x)$, $A_2'(\cdot\mid x)=\mathrm{Bern}(p_x+\varepsilon)$
with $p_1=1-\eta$, $p_0=\eta$. Every per-transition defect equals
$\varepsilon$, so $\eps{ctrl}=\varepsilon$ and the coarse bound
\eqref{eq:linear-bound} reads $\epsbar{ctrl}\le\Tc\,\eps{ctrl}=2\varepsilon$.

\paragraph{What the two definitions compute.}
With $r_t=A_t'/A_t$ and
$D(u_1,x):=\sum_{u_2}A_2(u_2\mid x)\,\lvert 1-r_1(u_1)\,r_2(u_2\mid x)\rvert$,
uniformity of $A_1$ gives
$\epsbar{ctrl}^{\mathrm{prod}}=\max_{x}\tfrac14[D(1,x)+D(0,x)]$ for (P) and
$\Delta_{\mathrm{cl}}=\tfrac14[D(1,x{=}1)+D(0,x{=}0)]$ for (H), because in
the loop branch $u_1$ is evaluated at \emph{its own} state $x_1=u_1$.
Enumerating the four outcomes (verified exactly by brute force):
$D(1,1)=4\varepsilon(1-\eta)+4\varepsilon^2$, $D(0,1)=D(1,0)=2\varepsilon$,
$D(0,0)=4\varepsilon(1-\eta)-4\varepsilon^2$, hence
\[
  \epsbar{ctrl}^{\mathrm{prod}}=\tfrac32\varepsilon-\varepsilon\eta+\varepsilon^2
  \qquad\text{versus}\qquad
  \Delta_{\mathrm{cl}}=2\varepsilon(1-\eta)\;>\;\epsbar{ctrl}^{\mathrm{prod}}
  \quad\text{whenever }\eta+\varepsilon<\tfrac12
\]
(e.g.\ $\varepsilon=\eta=0.1$: $0.150$ versus $0.180$, linear bound $0.20$;
$\varepsilon=\eta=0.01$: $0.0150$ versus $0.0198$, bound $0.02$).

\paragraph{Why the closed loop exceeds the product defect.}
$\Delta_{\mathrm{cl}}$ is a \emph{mixture of branch-wise maxima},
$\tfrac14\sum_{u_1}\max_x D(u_1,x)$, whereas (P) is the \emph{maximum of the
mixture}, $\max_x\tfrac14\sum_{u_1}D(u_1,x)$; the former dominates, strictly
when the worst-case conditioning differs across branches --- here the step-2
defect that aligns with the step-1 likelihood ratio sits at $x{=}1$ for
$u_1{=}1$ and at $x{=}0$ for $u_1{=}0$, and the loop delivers exactly that
state to each branch. Holding the trajectory fixed in (P) discards the
input--state correlation that makes defects add up. As $\eta\to0$,
$\Delta_{\mathrm{cl}}\to2\varepsilon=\Tc\,\eps{ctrl}$: the coarse linear
bound is attained, consistent with Remark~\ref{rem:linear}. Part~0 therefore
defines $\epsbar{\bullet}$ as the closed-loop hybrid displacement (H), for
which Lemma~\ref{lem:eps} is a triangle inequality by construction, and treats
(P)-type quantities only as what the state-matched audits estimate per
transition.
